% section: intro
\section{Introduction}
\label{sec:intro}

\subsection{The problem with picking a benchmark you are already good at}

A student project that fine-tunes a model has to show that the training did something. That sounds
trivial and is not. If you choose a task where a modern off-the-shelf model already performs well,
you can spend a month training and finish with a result that is flat, or worse than where you
started, and nothing to report.

Chart question answering is such a task. A current 2-billion-parameter vision-language model scores
\result{79.1} on ChartQA with no task-specific training at all, and the best published improvement
over a comparable prompted model is $+3.8$ points, obtained on four H100 GPUs. That is a poor place
to stake a project's only claim.

Chart \emph{grounding} --- pointing at the region of the image the answer came from --- is not
saturated. The strongest published result on RefChartQA's human subset is \result{32.83} AP@0.5,
produced by a \emph{larger} model trained for a single epoch with boxes emitted as plain text, an
objective carrying no spatial signal at all.

This work therefore reports both, measured against the same untrained model, and is arranged so
that the defensible measured improvement does not depend on the saturated one.

\subsection{Why chart grounding is hard, mechanically}

There is a concrete reason grounding scores are low, and it is not a matter of insufficient training.

A vision-language model does not see pixels. It divides the image into a grid and represents each
cell as one visual token. For the model used here that cell is $32 \times 32$ pixels. A target
narrower than one cell cannot be represented as a distinct thing --- it is blended with whatever
shares its cell.

Measured on RefChartQA validation data at the input resolution this system uses, \result{53.2\%} of
grounding targets are narrower than one visual token on at least one axis. The model is being asked
to point precisely at things smaller than the smallest unit it can perceive. Aggregate grounding
scores hide this completely, which is why every grounding result here is stratified by target size.

\subsection{What this system does}

Given a chart and a question, one model emits a single strict JSON record: whether the question is
answerable, the evidence regions with their labels and values, a typed expression tree, and a direct
answer. A small deterministic interpreter then validates the record and \emph{recomputes} the answer
from the expression tree.

The point of the separation is measurement, not accuracy. When the arithmetic is performed by an
interpreter rather than by the model, ``did it read the chart correctly'' and ``can it do
arithmetic'' become separately observable. Section~\ref{sec:analysis} uses that to attribute error
to seeing, to reasoning, or to execution.

\subsection{Contributions}

\begin{enumerate}
  \item A dual-target evaluation of one small vision-language model on ChartQA answer accuracy and
        RefChartQA grounding, each against a matched untrained baseline, using the official
        evaluators unmodified.
  \item \textbf{The first published zero-shot RefChartQA grounding measurement for this model.} No
        such number exists in the literature; producing it is a deliverable in its own right,
        whatever its value turns out to be.
  \item A measurement of what structured output costs --- the same weights evaluated with the JSON
        record required and not required --- on a model smaller than the two published estimates,
        asked for more than either.
  \item An error decomposition separating visual, reasoning and execution failure, made possible by
        the model/interpreter split.
  \item A set of negative and methodological findings, reported because they were expensive to
        obtain and are invisible from the published record: what the official grounding evaluator
        actually rewards (Section~\ref{sec:setup:metrics}), and the sub-token analysis above.
\end{enumerate}

\subsection{What this work does not claim}

Five statements are routinely conflated in this literature and are kept separate throughout:
official training reproduced; official checkpoint evaluated; matched student baseline trained;
before/after improvement achieved; published baseline exceeded. This work makes the third and
fourth. Section~\ref{sec:results} states each individually, including the ones that are negative.
